% Options for packages loaded elsewhere
\PassOptionsToPackage{unicode}{hyperref}
\PassOptionsToPackage{hyphens}{url}
\PassOptionsToPackage{dvipsnames,svgnames,x11names}{xcolor}
%
\documentclass[
  11pt,
]{article}
\usepackage{amsmath,amssymb}
\usepackage{iftex}
\ifPDFTeX
  \usepackage[T1]{fontenc}
  \usepackage[utf8]{inputenc}
  \usepackage{textcomp} % provide euro and other symbols
\else % if luatex or xetex
  \usepackage{unicode-math} % this also loads fontspec
  \defaultfontfeatures{Scale=MatchLowercase}
  \defaultfontfeatures[\rmfamily]{Ligatures=TeX,Scale=1}
\fi
\usepackage{lmodern}
\ifPDFTeX\else
  % xetex/luatex font selection
  \setmainfont[]{Latin Modern Roman}
  \setmonofont[]{Latin Modern Mono}
\fi
% Use upquote if available, for straight quotes in verbatim environments
\IfFileExists{upquote.sty}{\usepackage{upquote}}{}
\IfFileExists{microtype.sty}{% use microtype if available
  \usepackage[]{microtype}
  \UseMicrotypeSet[protrusion]{basicmath} % disable protrusion for tt fonts
}{}
\makeatletter
\@ifundefined{KOMAClassName}{% if non-KOMA class
  \IfFileExists{parskip.sty}{%
    \usepackage{parskip}
  }{% else
    \setlength{\parindent}{0pt}
    \setlength{\parskip}{6pt plus 2pt minus 1pt}}
}{% if KOMA class
  \KOMAoptions{parskip=half}}
\makeatother
\usepackage{xcolor}
\usepackage[margin=0.85in]{geometry}
\usepackage{longtable,booktabs,array}
\usepackage{calc} % for calculating minipage widths
% Correct order of tables after \paragraph or \subparagraph
\usepackage{etoolbox}
\makeatletter
\patchcmd\longtable{\par}{\if@noskipsec\mbox{}\fi\par}{}{}
\makeatother
% Allow footnotes in longtable head/foot
\IfFileExists{footnotehyper.sty}{\usepackage{footnotehyper}}{\usepackage{footnote}}
\makesavenoteenv{longtable}
\setlength{\emergencystretch}{3em} % prevent overfull lines
\providecommand{\tightlist}{%
  \setlength{\itemsep}{0pt}\setlength{\parskip}{0pt}}
\setcounter{secnumdepth}{-\maxdimen} % remove section numbering
\ifLuaTeX
\usepackage[bidi=basic]{babel}
\else
\usepackage[bidi=default]{babel}
\fi
\babelprovide[main,import]{american}
\ifPDFTeX
\else
\babelfont{rm}[]{Latin Modern Roman}
\fi
% get rid of language-specific shorthands (see #6817):
\let\LanguageShortHands\languageshorthands
\def\languageshorthands#1{}
\usepackage{amsmath,amssymb}
\usepackage{microtype}
\usepackage{fancyhdr}
\pagestyle{fancy}
\fancyhf{}
\renewcommand{\headrulewidth}{0pt}
\fancyfoot[C]{\thepage}
\setlength{\headheight}{14pt}
\setlength{\emergencystretch}{2em}
\setlength{\parskip}{4pt}
\setlength{\parindent}{0pt}
\AtBeginEnvironment{longtable}{\small}
\ifLuaTeX
  \usepackage{selnolig}  % disable illegal ligatures
\fi
\usepackage{bookmark}
\IfFileExists{xurl.sty}{\usepackage{xurl}}{} % add URL line breaks if available
\urlstyle{same}
\hypersetup{
  pdftitle={Local ordinary discriminants of Rogers Szegő polynomials at odd roots of unity},
  pdflang={en-US},
  colorlinks=true,
  linkcolor={black},
  filecolor={Maroon},
  citecolor={Blue},
  urlcolor={black},
  pdfcreator={LaTeX via pandoc}}

\title{Local ordinary discriminants of Rogers Szegő polynomials at odd
roots of unity}
\usepackage{etoolbox}
\makeatletter
\providecommand{\subtitle}[1]{% add subtitle to \maketitle
  \apptocmd{\@title}{\par {\large #1 \par}}{}{}
}
\makeatother
\subtitle{An all order derivation with separated and simple overlap
hypotheses}
\author{}
\date{Revised working mathematical note\\Version 2 \quad 11 October 2026}

\begin{document}
\maketitle

\section{1 Purpose and scope}\label{purpose-and-scope}

We derive the first moving-root terms of Rogers--Szegő polynomials at a
fixed odd root of unity, for every positive block multiplicity \(A\).
The finite \(q\)-binomial theorem gives an exact diagonal-operator
identity. A block-divisibility lemma then controls every order of the
formal remainder. In the case of one simple shared root, only one
additional product derivative is needed; its coefficient reduces to the
classical value at \(z=-1\).

Under the stated simplicity and separation hypotheses, these initial
forms determine the exact cyclotomic order and the complete leading
coefficient of the ordinary discriminant. We retain all signs and
distinguish the parameters \(q-\omega\), \(q^d-1\), and \(\Phi_d(q)\).

The all-order proof underlying this note received an independent scoped
acceptance on 11 October 2026. This note reorganizes that accepted proof
and its qualifications into a single mathematical account. Historical
publication priority remains unresolved. The conclusion is an explicit
all-order derivation from classical inputs; it does not assert that the
displayed formulas were previously published or establish a worldwide
originality claim. Source identities and the finite-control record
appear in the appendices.

Version 2 adds Appendix C, a concise integration of the independently
accepted finite-parameter analytic remainder, root-separation, and
discriminant-error certificate. It supplies explicit inputs to the
existing Q08 inverse-Vandermonde bound. The formal theorems and their
proofs are unchanged. Acceptance of the source mathematics is distinct
from independent document-fidelity review of this revision; Appendix D
identifies the frozen sources.

\section{2 Definitions and
conventions}\label{definitions-and-conventions}

Let \(K\) be a characteristic-zero field containing a primitive \(d\)th
root of unity \(\omega\), where \(d\geq 3\) is odd. Fix integers
\(A\geq 1\) and \(0\leq b<d\), and put \(n=Ad+b\). Algebraic closures
and formal power-series extensions of \(K\) may be used to discuss
individual roots.

The Gaussian binomial coefficient is the polynomial in \(\mathbb Z[q]\)
represented, for \(0\leq k\leq n\), by

\[
\genfrac{[}{]}{0pt}{}{n}{k}_q=\prod_{j=1}^{k}\frac{1-q^{n-k+j}}{1-q^j}.
\]

Specialization at a root of unity means evaluation of this polynomial,
rather than substitution into an unreduced quotient. Define

\[
H_n(z;q)=\sum_{k=0}^{n}\genfrac{[}{]}{0pt}{}{n}{k}_q z^k,
\qquad D_n(q)=\operatorname{Disc}_z H_n(z;q),
\]

\[
F(z)=1+z^d,\qquad G(z)=H_b(z;\omega),\qquad
P_A=\prod_{j=1}^{A}j^j.
\]

Every \(H_n\) is monic in \(z\). For a monic polynomial
\(f(z)=\prod_i(z-r_i)\), the ordinary discriminant and the resultant
with a monic polynomial \(g\) are

\[
\operatorname{Disc}(f)=\prod_{i<j}(r_i-r_j)^2,
\qquad \operatorname{Res}(f,g)=\prod_i g(r_i).
\]

We use
\(\operatorname{Disc}(1)=\operatorname{Disc}(\text{monic linear})=1\)
and empty products equal to \(1\). All discriminants below use the
ordinary \(z\) derivative, not a \(q\)-difference operator.

The monic probabilists' Hermite polynomial is

\[
\operatorname{He}_m(y)=m!\sum_{j=0}^{\lfloor m/2\rfloor}
\frac{(-1)^j y^{m-2j}}{2^j j!(m-2j)!}.
\]

In particular \(\operatorname{He}_0=1\), \(\operatorname{He}_1=y\), and
\(\operatorname{Disc}(\operatorname{He}_A)=P_A\). Write
\((a;q)_m=\prod_{j=0}^{m-1}(1-aq^j)\).

We use formal local parameters

\[
q=\omega e^h,\qquad q(\varepsilon)=\omega(1+\varepsilon^2)^{1/d},
\qquad h=\frac{\log(1+\varepsilon^2)}{d}.
\]

The binomial and logarithmic series have their constant terms fixed as
written. Thus \(q(0)=\omega\), \(q(\varepsilon)^d-1=\varepsilon^2\), and
\(h=\varepsilon^2/d+O(\varepsilon^4)\). A remainder \(O(\varepsilon^s)\)
in a local polynomial identity is divisible by \(\varepsilon^s\) in
\(K[x][[\varepsilon]]\), after adjoining square roots if convenient.
These square roots cancel from the displayed initial forms
coefficientwise.

\section{3 The two local theorems}\label{the-two-local-theorems}

\subsection{3 1 Separated roots}\label{separated-roots}

\textbf{Theorem 1.} For every root \(\alpha\) of \(F\), the following
coefficient identity holds:

\[
H_n(\alpha+\varepsilon x;q(\varepsilon))
=\varepsilon^A(-1)^A d^{A/2}G(\alpha)
\operatorname{He}_A\!\left(\frac{\sqrt d\,x}{\alpha}\right)
+O(\varepsilon^{A+1}). \tag{1}
\]

If \(G(\alpha)=0\), the written coefficient vanishes and (1) does not
identify the first nonzero term. For the separated discriminant
conclusion assume

\[
\operatorname{Disc}(G)\operatorname{Res}(F,G)\ne 0. \tag{2}
\]

Put \(M_{\mathrm{sep}}=dA(A-1)/2\). Then

\[
D_n(q(\varepsilon))=
\varepsilon^{2M_{\mathrm{sep}}}
\bigl(C_{\mathrm{sep}}+O(\varepsilon^2)\bigr), \tag{3}
\]

\[
C_{\mathrm{sep}}=
(-1)^{A(d-1)/2}d^{dA(A+1)/2}P_A^d
\operatorname{Disc}(G)\operatorname{Res}(F,G)^{2A}. \tag{4}
\]

The coefficient is nonzero, and
\(\operatorname{ord}_{\Phi_d}D_n=M_{\mathrm{sep}}\). The same separated
proof also covers the classical endpoint \(d=1\), \(b=0\); this endpoint
is not presented as a new result.

\subsection{3 2 One simple shared root}\label{one-simple-shared-root}

Suppose \(b\) is odd, so \(1\leq b<d\), and write

\[
G(z)=(z+1)U(z),\qquad g=U(-1)\ne 0.
\]

Set

\[
m=\frac{d+b}{2}<d,\qquad
W=\frac{(\omega;\omega)_b}{(\omega^2;\omega^2)_m},
\qquad \tau=\frac{dW}{g}, \tag{5}
\]

\[
J_{A+1}(y;\tau)=\operatorname{He}_{A+1}(y)
+A(1-\tau)\operatorname{He}_{A-1}(y). \tag{6}
\]

Every denominator factor in \(W\) is nonzero: since \(d\) is odd,
\(\omega^2\) also has order \(d\), and \(m<d\).

\textbf{Theorem 2.} Under \(g\ne 0\), the shared-root coefficient
identity is

\[
H_n(-1+\varepsilon x;q(\varepsilon))
=\varepsilon^{A+1}d^{(A-1)/2}g
J_{A+1}(\sqrt d\,x;\tau)+O(\varepsilon^{A+2}). \tag{7}
\]

Both Hermite summands have parity \(A+1\). Therefore the leading factor
in (7) is equivalently
\((-1)^{A+1}d^{(A-1)/2}gJ_{A+1}(-\sqrt d\,x;\tau)\), as in the accepted
source statement.

For the global overlap discriminant additionally assume

\[
\operatorname{Disc}(U)\operatorname{Res}(F,U)
\operatorname{Disc}(J_{A+1})\ne 0. \tag{8}
\]

Here \(\operatorname{Disc}(J_{A+1})\) means the discriminant in \(y\) at
the specified value of \(\tau\). These hypotheses give exactly one
shared residual root, namely \(-1\), simple in \(G\), with distinct
first displacements. In fact \(\operatorname{Res}(F,U)\ne0\) already
implies \(g\ne0\).

Put \(M_{\mathrm{ov}}=dA(A-1)/2+A\). Then

\[
D_n(q(\varepsilon))=
\varepsilon^{2M_{\mathrm{ov}}}
\bigl(C_{\mathrm{ov}}+O(\varepsilon^2)\bigr), \tag{9}
\]

\[
\begin{aligned}
C_{\mathrm{ov}}={}&(-1)^{A(d-1)/2}
 d^{dA(A+1)/2+A}P_A^{d-1}\operatorname{Disc}(J_{A+1})\\
&\mathrel{}\cdot\operatorname{Disc}(U)
\operatorname{Res}(F,U)^{2A}g^2.
\end{aligned} \tag{10}
\]

The coefficient is nonzero, and
\(\operatorname{ord}_{\Phi_d}D_n=M_{\mathrm{ov}}\). No assertion removes
(2) or (8), extends this proof to even \(d\), or gives a
positive-characteristic unit conclusion.

\section{4 The exact operator and the all order
bound}\label{the-exact-operator-and-the-all-order-bound}

The classical finite \(q\)-binomial identity is

\[
\mathcal P_n(z;q)=\prod_{j=0}^{n-1}(1+q^jz)
=\sum_{k=0}^{n}q^{k(k-1)/2}\genfrac{[}{]}{0pt}{}{n}{k}_qz^k. \tag{11}
\]

It is recorded, for example, in Gorodetsky {[}1, equation (1.5){]}.
Define diagonal and differential operators on polynomials by

\[
T_\omega(z^k)=\omega^{k(k-1)/2}z^k,
\qquad \mathcal L=z^2\partial_z^2,
\]

and set \(R_n(z,h)=T_\omega^{-1}\mathcal P_n(z;\omega e^h)\). Comparing
coefficients in (11) gives the exact formal identity

\[
H_n(z;\omega e^h)=\exp(-h\mathcal L/2)R_n(z,h). \tag{12}
\]

This takes place in the \(h\)-adic module of polynomials in \(z\) of
degree at most \(n\). The coefficient of a fixed \(h\) power uses only
finitely many terms of the exponential.

Oddness of \(d\) is essential to the next step. The multiplier ratio for
\(z^{k+d}\) and \(z^k\) is

\[
\omega^{(k+d)(k+d-1)/2-k(k-1)/2}
=\omega^{dk+d(d-1)/2}=1.
\]

Consequently \(T_\omega\) and its inverse commute with multiplication by
\(z^d\) and by every power of \(F\). For even \(d\) the same ratio is
\(-1\).

\textbf{Lemma 3.} If \(R_n(z,h)=\sum_{r\geq0}h^rR_{n,r}(z)\), then

\[
F^{\max(A-r,0)}\mid R_{n,r}\quad\text{for every }r\geq0,
\qquad R_{n,0}=F^AG. \tag{13}
\]

\textbf{Proof.} Partition the first \(Ad\) factors of \(\mathcal P_n\)
into \(A\) consecutive blocks of length \(d\), followed by a residual
block of length \(b\). At \(h=0\), each full block equals \(F\). In a
term of the \(r\)th product derivative, at most \(r\) full blocks are
differentiated. At least \(A-r\) intact blocks remain if \(r\leq A\).
Apply \(T_\omega^{-1}\), which commutes with their product. Division by
\(r!\) does not affect divisibility. The residual finite \(q\)-binomial
identity gives \(R_{n,0}=F^AG\). This proves (13) at every order.
\(\square\)

At a root \(\alpha\) of \(F\), put \(u=z-\alpha\) and assign \(u\)
weight \(1\) and \(h\) weight \(2\). The \(r\)th summand of \(R_n\) has
weighted order at least

\[
2r+\max(A-r,0). \tag{14}
\]

This is at least \(A+1\) for \(r\geq1\) and at least \(A+2\) for
\(r\geq2\), including \(r>A\).

In these local coordinates,

\[
h\mathcal L=h\alpha^2\partial_u^2
+2h\alpha u\partial_u^2+hu^2\partial_u^2. \tag{15}
\]

The three terms preserve weight, increase it by one, and increase it by
two, respectively. Hence \(\exp(-h\mathcal L/2)\) preserves every
weighted remainder ideal. On a fixed leading weight its action is
\(\exp(-h\alpha^2\partial_u^2/2)\). Its weight-preserving part is a
finite sum on each weighted homogeneous polynomial, because repeated
second derivatives terminate. The full expression is also well defined
\(h\)-adically by (12). No analytic convergence is used.

\section{5 Proof of the initial forms}\label{proof-of-the-initial-forms}

\subsection{5 1 The separated
coefficient}\label{the-separated-coefficient}

Since \(F'(\alpha)=-d/\alpha\), the weight-\(A\) part of \(R_n\) is
\nopagebreak[4]
\[
(-d/\alpha)^A G(\alpha)u^A.
\]

All positive-\(r\) parts have larger weight by (14). The finite
polynomial identity

\[
\exp(-s\partial_u^2/2)u^A
=\sum_{j=0}^{\lfloor A/2\rfloor}
\frac{A!(-s)^j u^{A-2j}}{2^j j!(A-2j)!} \tag{16}
\]

equals \(s^{A/2}\operatorname{He}_A(u/\sqrt s)\). Apply (16) with
\(s=h\alpha^2\), then put \(u=\varepsilon x\) and
\(h=\varepsilon^2/d+O(\varepsilon^4)\). The result is (1). Formula (16)
also explains directly why square roots cancel.

\subsection{5 2 The first product derivative at the shared
root}\label{the-first-product-derivative-at-the-shared-root}

Let \(\mathcal D=z\partial_z\) and define

\[
E_b(z)=\prod_{r=0}^{b-1}(1+\omega^rz),
\]

\[
S(z)=\sum_{r=0}^{d-1}r\omega^rz
\prod_{\substack{0\leq s<d\\s\ne r}}(1+\omega^sz),
\]

with \(S_b\) defined by the same sum and product over \(0\leq r,s<b\).
The derivative of full block \(i\), indexed from \(i=0\), is
\(S+id^2z^d\). Therefore

\[
\begin{aligned}
\left.\partial_h\mathcal P_n(z;\omega e^h)\right|_{h=0}
={}&F^{A-1}\left(AE_bS+\frac{d^2A(A-1)}2z^dE_b\right)\\
&+F^A\left(Ad\mathcal D E_b+S_b\right).
\end{aligned} \tag{17}
\]

Apply \(T_\omega^{-1}\). The \(z^dE_b\) term acquires a factor \(G\),
which vanishes at \(-1\). The final \(F^A\) term has weight at least
\(A+2\) after multiplication by \(h\). Thus, with \(u=z+1\) and

\[
c_b=\bigl(T_\omega^{-1}(E_bS)\bigr)(-1),
\]

the only weight-\((A+1)\) contribution from \(hR_{n,1}\) is
\(Ad^{A-1}c_bhu^{A-1}\). This statement also holds for \(A=1\); the term
with factor \(A(A-1)\) is then zero.

We can determine \(c_b\) using \(A=1\). Differentiating (12) gives

\[
R_{d+b,1}=\left.\partial_hH_{d+b}(z;\omega e^h)\right|_{h=0}
+\frac{\mathcal L}{2}(FG).
\]

Because \(FG=dg\,u^2+O(u^3)\), the second term equals \(dg\) at
\(z=-1\). The classical Gaussian value is

\[
H_{2m}(-1;q)=(q;q^2)_m
=\frac{(q;q)_{d+b}}{(q^2;q^2)_m},\qquad 2m=d+b. \tag{18}
\]

Odd-degree values vanish. These identities also follow from the
three-term recurrence displayed in Appendix A; see Schultz--Walker {[}2,
printed page 2{]}. The numerator on the right of (18) contains exactly
one vanishing factor at \(q=\omega\), namely \(1-q^d\). The remaining
factors multiply to \(d(\omega;\omega)_b\), since
\(\prod_{j=1}^{d-1}(1-\omega^j)=d\). Also
\(\partial_h(1-e^{dh})|_{h=0}=-d\). Therefore

\[
H_{d+b}(-1;\omega e^h)=-d^2Wh+O(h^2),
\qquad c_b=dg-d^2W=dg(1-\tau). \tag{19}
\]

The weight-\((A+1)\) static term is \(d^Agu^{A+1}\), while every
\(r\geq2\) term has weight at least \(A+2\). The complete initial form
of \(R_n\) is consequently

\[
d^Ag\left(u^{A+1}+A(1-\tau)hu^{A-1}\right). \tag{20}
\]

At \(\alpha=-1\), apply \(\exp(-h\partial_u^2/2)\) to (20) and
substitute \(u=\varepsilon x\), \(h=\varepsilon^2/d+O(\varepsilon^4)\).
Formula (16) gives (7). Both the factor \(d\) in \(\tau=dW/g\) and the
factor \(1/d\) in the leading term of \(h\) are necessary.

\section{6 Ordinary discriminants and
signs}\label{ordinary-discriminants-and-signs}

Under the respective guards, each simple root of the initial polynomial
lifts by the formal implicit-function theorem to a root branch in an
algebraic closure with coefficients in formal \(\varepsilon\) series. In
the separated case there are \(A\) branches at each of the \(d\) roots
of \(F\), together with \(b\) residual branches. In the overlap case
there are \(A+1\) branches at \(-1\), \(A\) at each other root of \(F\),
and \(b-1\) residual branches. These counts account for all \(n\) roots.

\subsection{6 1 Separated clusters}\label{separated-clusters}

If \(h_1,\ldots,h_A\) are the roots of \(\operatorname{He}_A\), the
first displacements at \(\alpha\) are \(\alpha h_j/\sqrt d\). The
product of squared differences within all clusters is therefore, to
leading order,

\[
\varepsilon^{dA(A-1)}d^{-dA(A-1)/2}P_A^d. \tag{21}
\]

The omitted \(\alpha\) powers cancel: \(\prod_{F(\alpha)=0}\alpha=-1\)
and \(A(A-1)\) is even. The cross-cluster product is
\(\operatorname{Disc}(F)^{A^2}\), where

\[
\operatorname{Disc}(F)=(-1)^{(d-1)/2}d^d. \tag{22}
\]

Residual pairs and pairs between clusters and residual roots contribute
\(\operatorname{Disc}(G)\operatorname{Res}(F,G)^{2A}\). Multiplying
these terms, using \(A^2\equiv A\pmod2\), proves (3) and (4).

\subsection{6 2 The shared cluster}\label{the-shared-cluster}

The first displacements in (7) are \(j_k/\sqrt d\), where the \(j_k\)
are the roots of \(J_{A+1}\). The parity convention with negative
argument gives the same set. Put \(F_0=F/(z+1)\). Then

\[
F_0(-1)=d,\qquad
\operatorname{Disc}(F_0)=(-1)^{(d-1)/2}d^{d-2}.
\]

The product of the roots of \(F_0\) is \(1\). The within-cluster
contribution is

\[
\varepsilon^{2M_{\mathrm{ov}}}d^{-M_{\mathrm{ov}}}
P_A^{d-1}\operatorname{Disc}(J_{A+1}). \tag{23}
\]

Cross-cluster pairs contribute
\nopagebreak[4]
\[
\operatorname{Disc}(F_0)^{A^2}d^{2A(A+1)}
=(-1)^{A(d-1)/2}d^{dA^2+2A}. \tag{24}
\]

The residual contribution is
\(\operatorname{Disc}(U)\operatorname{Res}(F,U)^{2A}g^2\). The extra
factor \(g^2\) accounts for the extra root in the shared cluster.
Combining (23) and (24) proves (9) and (10).

In both cases, nonvanishing of the displayed coefficient establishes the
exact order. If a guard fails, these root-lifting arguments do not
certify that order. Although individual root branches need not be even
in \(\varepsilon\), the discriminant is even because \(q(\varepsilon)\)
is even. This gives the normalized remainders \(O(\varepsilon^2)\) in
(3) and (9).

\section{7 Conversion between local
parameters}\label{conversion-between-local-parameters}

Let \((M,C)\) mean either \((M_{\mathrm{sep}},C_{\mathrm{sep}})\) or
\((M_{\mathrm{ov}},C_{\mathrm{ov}})\), under its hypotheses, and put
\(t=q^d-1\). The formulas above say

\[
D_n(q)=t^M\bigl(C+O(t)\bigr)\quad\text{locally at }q=\omega. \tag{25}
\]

Since \(t=d\omega^{d-1}(q-\omega)+O((q-\omega)^2)\), the coefficient of
\((q-\omega)^M\) is

\[
(d\omega^{d-1})^M C. \tag{26}
\]

It is generally different from \(C\). Let \(Q_d(q)=(q^d-1)/\Phi_d(q)\).
The integral polynomial \(D_n(q)\) has the same exact root multiplicity
at every primitive \(d\)th root by Galois conjugacy. Monic division
gives

\[
D_n(q)=\Phi_d(q)^M E(q),\qquad E(q)\in\mathbb Z[q],
\qquad E(\omega)=Q_d(\omega)^M C. \tag{27}
\]

These assertions are in characteristic zero. Nonvanishing of \(C\) there
does not imply that its reduction is a unit in any specified positive
characteristic.

\section{8 Classical antecedents and logical
dependence}\label{classical-antecedents-and-logical-dependence}

The input (11) is classical. The static identity
\(H_n(z;\omega)=F(z)^AG(z)\) is the usual root-of-unity factorization;
(12)--(13) also derive it directly. The special value (18), the Hermite
discriminant, and the root-product discriminant rule are established
ingredients. The proof here supplies the block filtration and the first
surviving coefficients explicitly.

Schultz--Walker {[}2, Theorem 1.1, printed page 3{]} provides a relevant
all-order vanishing antecedent. If \(\alpha^d=-1\) has order \(2c\),
then \(c\mid d\) and \(d/c\) is odd. In their theorem take \(k=d/c\),
their auxiliary exponent and congruence-order parameter equal to zero,
derivative index \(r\), and weight polynomial
\(P(s)=s(s-1)\cdots(s-\ell+1)\). Up to the nonzero factor
\(\alpha^{-\ell}\), their sum is
\(\partial_q^r\partial_z^\ell H_n(\alpha;q)\). Their bound becomes
\(n\geq(\ell+2r+1)d\). Hence

\[
\partial_q^r\partial_z^\ell H_n(\alpha;\omega)=0
\quad\text{whenever }2r+\ell<A. \tag{28}
\]

Their primitive-choice passage is on printed page 13. The stated
implication covers roots of smaller order as well as primitive \(2d\)th
roots. The cited theorem gives vanishing; it does not specify the
surviving coefficient.

The logical distinctions matter. For \(b=0\) and \(A\geq2\), adding
\(\lambda(q^d-1)F^{A-2}\) to a monic polynomial family preserves its
static specialization and all below-weight-\(A\) vanishings, but changes
its weight-\(A\) initial form. At the shared root, adding
\(\lambda(q^d-1)F^{A-1}\) similarly changes weight \(A+1\) while
preserving the lower vanishings. These are logical counterexamples to
determining a coefficient from vanishing data alone. The perturbed
families are not asserted to be Rogers--Szegő polynomials.

Conversely, a full local initial-form identity implies its lower-weight
mixed-derivative vanishings. The coordinate change
\(q-\omega=\omega h+O(h^2)\) is invertible and filtered; substitution
\(u=\varepsilon x\), \(h=\varepsilon^2/d+O(\varepsilon^4)\) separates
the coefficients of each weight by powers of \(x\). This local converse
does not recover the full Schultz--Walker theorem. Likewise, the guarded
initial forms imply the discriminant coefficients, but a single
discriminant coefficient does not determine all initial polynomials or
imply the operator identity.

The accepted implication is therefore

\[
\begin{gathered}
\text{finite product identity and proved block filtration}\\
+\ \text{shared root special value when needed}\\
\Longrightarrow\ \text{local initial forms}\\
\Longrightarrow\ \text{guarded discriminant formulas}.
\end{gathered}
\]

This derivation settles the previously unworked all-\(A\) implication
route. It does not settle when the resulting formulas were first
published. The finite controls in Appendix A support the audit and
cannot establish the all-\(A\) claim by themselves.

\section{9 Remaining comparisons and
extensions}\label{remaining-comparisons-and-extensions}

Historical comparison still requires the exact six-page body of
Zhedanov's 1997 paper {[}4{]}, the printed page 360 image of
Lubinsky--Saff {[}3{]}, and substantive bodies along the relevant
Lubinsky--Saff citation chain. Missing or inaccessible text supplies no
evidence that a theorem is absent. The comparison of the
complex-parameter disk and inverse-pair applications is a separate open
priority question.

The local theorem itself retains several boundaries. Multiple common
residual roots, a repeated root at \(-1\), or a repeated root of
\(J_{A+1}\) would require a different initial-form analysis. The proof
does not remove the composite residual squarefreeness or resultant
guards. Even \(d\) changes the periodicity used in (13). Uniform
estimates when degree grows, and arithmetic conclusions requiring
positive-characteristic units, need additional arguments. None of these
extensions is claimed here. Appendix C supplies conservative bounds for
each specified finite instance; it does not supply efficient or uniform
growing-degree estimates, a numerical implementation certificate, or
any change in historical-priority status.

For \(A=1\), the separated discriminant is already the static
product-discriminant consequence; the simple-overlap calculation reduces
to (18) and a Taylor expansion. For arbitrary \(A\), the accepted result
is a fully written all-order coefficient derivation under the guards.
Any assessment of its publication contribution should separate
correctness, explicit derivability, historical priority, and the value
of particular applications.

\section{References}\label{references}

{[}1{]} Ofir Gorodetsky. \emph{q-congruences, with applications to
supercongruences and the cyclic sieving phenomenon}. arXiv:1805.01254v2,
equation (1.5), preprint page 2. The citation records the classical
finite \(q\)-binomial identity and does not assign its origin to this
paper. \href{https://arxiv.org/pdf/1805.01254v2}{Primary preprint}

{[}2{]} Andrew Schultz and Robert Walker. \emph{A generalization of the
Gaussian formula and a q-analog of Fleck's congruence}.
arXiv:1202.0199v1, printed pages 2--3 and 13. Published in \emph{Journal
of Number Theory} 133(11) (2013), 3717--3738, DOI
10.1016/j.jnt.2013.05.003. Page locators here refer to the preprint.
\href{https://arxiv.org/pdf/1202.0199v1}{Primary preprint}

{[}3{]} D. S. Lubinsky and E. B. Saff. \emph{Convergence of Padé
approximants of partial theta functions and the Rogers--Szegő
polynomials} (1987). Historical comparison remains incomplete at the
source scope stated in Section 9.
\href{https://doi.org/10.1007/BF01890574}{Publisher record}

{[}4{]} A. Zhedanov. \emph{Gauss sums and orthogonal polynomials}
(1997). The exact body remains an outstanding comparison source.
\href{https://doi.org/10.1142/S0217751X97000438}{Publisher record}

\section{Appendix A Finite controls and
reproducibility}\label{appendix-a-finite-controls-and-reproducibility}

The independent review used a freshly written three-term recurrence
implementation:

\[
H_0=1,\quad H_1=1+z,\quad
H_n=(1+z)H_{n-1}+(q^{n-1}-1)zH_{n-2}.
\]

It computed exact rational truncated \(\varepsilon\) series in
\(\mathbb Q[w]/(w^d-1)\), then reduced modulo \(\Phi_d(w)\). It did not
generate the tested jets using the operator proof or the author's
\(q\)-Pascal program. The author's code was read as inert text during
that review, not imported or executed.

All 16 saved moving-root controls passed. In the table, \(s\) specifies
the tested root \(\alpha=-\omega^s\). ``Other root'' means evaluation of
formula (1); it does not claim that every residual guard holds globally
for that parameter triple.

\begin{longtable}[]{@{}rrrrlr@{}}
\toprule\noalign{}
\(A\) & \(d\) & \(b\) & \(s\) & Local formula & Initial weight \\
\midrule\noalign{}
\endhead
\bottomrule\noalign{}
\endlastfoot
1 & 3 & 1 & 0 & Shared root & 2 \\
2 & 3 & 2 & 2 & Other root & 2 \\
4 & 3 & 1 & 0 & Shared root & 5 \\
6 & 3 & 1 & 0 & Shared root & 7 \\
10 & 3 & 1 & 0 & Shared root & 11 \\
12 & 3 & 0 & 2 & Other root & 12 \\
5 & 5 & 1 & 0 & Shared root & 6 \\
4 & 5 & 3 & 0 & Shared root & 5 \\
3 & 5 & 4 & 1 & Other root & 3 \\
3 & 7 & 5 & 0 & Shared root & 4 \\
4 & 7 & 2 & 3 & Other root & 4 \\
3 & 9 & 7 & 0 & Shared root & 4 \\
4 & 9 & 2 & 3 & Other root & 4 \\
2 & 15 & 3 & 5 & Other root & 2 \\
1 & 15 & 13 & 0 & Shared root & 2 \\
7 & 1 & 0 & 0 & Classical endpoint & 7 \\
\end{longtable}

Five direct ordinary-discriminant computations also passed the exact
cyclotomic valuation and both coefficient conversions:

\begin{longtable}[]{@{}lrr@{}}
\toprule\noalign{}
\((A,d,b)\) & Degree \(n\) & Exact \(\Phi_d\) order \\
\midrule\noalign{}
\endhead
\bottomrule\noalign{}
\endlastfoot
\((1,3,2)\) & 5 & 0 \\
\((1,5,1)\) & 6 & 1 \\
\((2,3,2)\) & 8 & 3 \\
\((2,3,1)\) & 7 & 5 \\
\((1,5,3)\) & 8 & 1 \\
\end{longtable}

Six mutation controls detected the intended errors: omitting the heat
term; replacing \(\tau\) by \(1\); omitting \(d\) from \(\tau\);
assuming the same periodicity for even \(d\); omitting the \(q-\omega\)
scaling; and dropping the separated resultant guard. Constant-polynomial
discriminants were explicitly set to \(1\).

This appendix reports the sealed independent-control run. The full exact
outputs, executable source, and run log are retained in the independent
review archive identified below. They test bounded examples; the proof
for every \(A\) is Lemma 3 and Sections 5--7.

\section{Appendix B Source identities and review
boundary}\label{appendix-b-source-identities-and-review-boundary}

The original proof archive and independent review archive retain their
existing identities. This working note is a separate document, not a
replacement or a retrospective edit of either archive. The prior
independent acceptance concerns the exact proof bytes below; editorial
consistency of this note must be distinguished from a new independent
mathematical audit.

\textbf{All order implication archive}

F020\_All\_Order\_Implication\_Provisional.zip

Library identity: libfile\_f90f326d1b0c8191a838d4cb2e7fab76, version 0.

SHA256:
\texttt{79273ecedccd68bf0d174c843eb5cecc9d32cbb596b115519d26da11352b9ffd}

Accepted proof member ALL\_ORDER\_IMPLICATION.md

SHA256:
\texttt{8ca570c7e7c41c50827b1dea3feb931cbcfe0fdab3c3a2016c6e7d0989126abc}

\textbf{Independent review archive}

F020\_All\_Order\_Independent\_Review.zip

Library identity: libfile\_770d307623188191a6408f4915be0687, version 0.

SHA256:
\texttt{45671b26ad19d741ca9a8b08da11db23d946a60f6758f18b98544b2a1a8f860e}

Independent review member INDEPENDENT\_REVIEW.md

SHA256:
\texttt{aa6fab3492f265fca7a667a213f57f1453dcfb280d94280ce74b330b353bffd3}

Independent control output INDEPENDENT\_CONTROLS.json

SHA256:
\texttt{778d7585963a6bef193280895da256cc6cac6ee6c6b5ab95e8a2f8ac0216f0ab}

Independent control source independent\_controls.py

SHA256:
\texttt{2ddd2cbef7c652152aae23f7877494bec320ff6c4e8ec4e9c56cd2847cca14a0}

The independent verdict accepts the characteristic-zero all-\(A\)
implication in the separated and one-simple-overlap scope, with no
correction required to its displayed formulas. It checked selected
primary PDF text for {[}1{]} and {[}2{]}. Its attempted primary
page-image retrieval failed; it does not claim a rendered-page
inspection or an exhaustive whole-paper negative result.

For the earlier separated audit, the recovered file has the following
SHA256.

\texttt{f093175f18d1941a99a605a5d816b1db31e5b2c4161ee7df4bf29291ff8eff3c}

The catalogue lists a distinct export SHA256.

\texttt{4b86b3078bbdf644a3317e67449a974c356e8b2dd8a9f66a393df71341d3e8d9}

It is not silently identified with those bytes. The original separated
and overlap reports, and their independent acceptance records, remain
unchanged.

\clearpage
\section{Appendix C Finite parameter analytic certificates}
\label{appendix-c-finite-certificates}

This appendix turns the accepted initial forms into explicit analytic
error and separation bounds for fixed finite parameters. Its source
certificate received an independent mathematical acceptance on
11 October 2026. The argument is a conservative application of Cauchy's
estimate and Rouch\'e's theorem. It makes no claim of optimal constants,
efficient growing-degree behavior, implementation stability, or historical
novelty. Appendix D binds this integration to the accepted source bytes.

\subsection{C 1 Analytic branch and normalized remainder}

Fix a complex embedding of the data in Section 2, with odd \(d\geq3\),
\(A\geq1\), and \(0\leq b<d\). Retain either the separated guard (2)
or all the simple-overlap hypotheses of Section 3.2, including odd \(b\)
and (8). On \(|\varepsilon|<1\), use
\[
q(\varepsilon)=\omega\exp\!\left(\frac{\operatorname{Log}(1+\varepsilon^2)}d\right).
\]
The principal logarithm is analytic here because
\(\operatorname{Re}(1+\varepsilon^2)>0\). Its Taylor series agrees with
the formal branch in Section 2. Every root \(\alpha\) of \(F\) has
\(|\alpha|=1\). Write its initial form as
\[
H_n(\alpha+\varepsilon x;q(\varepsilon))
=\varepsilon^{m_\alpha}\bigl(P_\alpha(x)+O(\varepsilon)\bigr).
\tag{C1}
\]
At an ordinary center, \(m_\alpha=A\) and \(P_\alpha\) is the
coefficient in (1); at the exceptional center \(-1\),
\(m_{-1}=A+1\) and \(P_{-1}\) is the coefficient in (7).
If \(c_\alpha\) denotes the leading coefficient, then
\[
|c_\alpha|=d^A|G(\alpha)|\quad\text{at an ordinary center},
\qquad |c_{-1}|=d^A|g|\quad\text{at the exceptional center}.
\]
All these \(P_\alpha\) have simple roots under the retained guards.

For \(Q\geq1\) and \(Z\geq0\), define the positive majorant
\[
\mathcal B_n(Q,Z)=\sum_{k=0}^n\binom nk Q^{k(n-k)}Z^k
\leq Q^{\lfloor n^2/4\rfloor}(1+Z)^n,
\qquad Q_s=(1+s^2)^{1/d}. \tag{C2}
\]
For \(X\geq0\), \(0<r<s<1\), \(0<|\varepsilon|\leq r\), and
\(|x|\leq X\), the explicit normalized remainder is
\[
\left|\varepsilon^{-m_\alpha}
 H_n(\alpha+\varepsilon x;q(\varepsilon))-P_\alpha(x)\right|
\leq E_\alpha(r;s,X)
:=\frac{\mathcal B_n(Q_s,1+sX)}{s^{m_\alpha}}\frac r{s-r}.
\tag{C3}
\]
Indeed, the Gaussian polynomial has nonnegative integer coefficients,
coefficient sum \(\binom nk\), and degree \(k(n-k)\), by its
polynomial Pascal recurrence. On \(|\varepsilon|=s\), its contribution
is bounded by (C2). Cauchy's estimate bounds the \(\ell\)th coefficient
of the unnormalized function by
\(\mathcal B_n(Q_s,1+sX)s^{-\ell}\). Equation (C1) cancels the
coefficients below \(m_\alpha\); summing the tail proves (C3).
In particular the normalization factor \(s^{-m_\alpha}\) is essential.

\subsection{C 2 Certified moving root tubes}

Write \(P_\alpha(x)=c_\alpha\prod_{j=1}^{m_\alpha}(x-\xi_{\alpha,j})\).
For \(m_\alpha>1\), let \(\delta_\alpha\) and \(D_\alpha\) be the
minimum and maximum pairwise direction distances, and choose
\(0<\eta_\alpha<\delta_\alpha/2\). For \(m_\alpha=1\), choose any
\(\eta_\alpha>0\). Set
\[
X_\alpha=\max_j|\xi_{\alpha,j}|+\eta_\alpha,\qquad
\mu_\alpha=|c_\alpha|\eta_\alpha
 \min_j\prod_{k\ne j}(|\xi_{\alpha,j}-\xi_{\alpha,k}|-\eta_\alpha).
\tag{C4}
\]
The empty product for a singleton is one. If
\[
E_\alpha(r;s,X_\alpha)<\mu_\alpha,
\quad\text{equivalently}\quad
r<\frac{\mu_\alpha s^{m_\alpha+1}}
 {\mathcal B_n(Q_s,1+sX_\alpha)+\mu_\alpha s^{m_\alpha}},
\tag{C5}
\]
then for every \(0<|\varepsilon|\leq r\), each disk
\[
|z-(\alpha+\varepsilon\xi_{\alpha,j})|<|\varepsilon|\eta_\alpha
\tag{C6}
\]
contains exactly one zero, counted with multiplicity, of \(H_n\).
It is therefore a simple zero. For two roots in this cluster,
\[
|\varepsilon|(\delta_\alpha-2\eta_\alpha)
< |z_i-z_j|
< |\varepsilon|(D_\alpha+2\eta_\alpha). \tag{C7}
\]
To prove this, the product in (C4) bounds \(|P_\alpha|\) from below
on each circle \(|x-\xi_{\alpha,j}|=\eta_\alpha\). Rouch\'e's theorem
and (C3) preserve the single-zero count. Rescaling by the nonzero
complex number \(\varepsilon\) proves (C6), and the triangle inequality
proves (C7). The bound in (C5) is positive and less than \(s\), so every
guarded finite instance admits a nonzero local radius.

An explicit quadratic root error is also available. Fix an admissible
\(\eta_0\), put \(X_0=\max_j|\xi_{\alpha,j}|+\eta_0\), and set
\[
L_\alpha=|c_\alpha|\min_j\prod_{k\ne j}
 (|\xi_{\alpha,j}-\xi_{\alpha,k}|-\eta_0).
\]
If \(2E_\alpha(r;s,X_0)/L_\alpha\leq\eta_0\), apply (C5) with
\(\eta_\alpha=2E_\alpha(r;s,X_0)/L_\alpha\): its boundary lower bound
is at least twice the error. Thus
\[
|z_j-(\alpha+\varepsilon\xi_{\alpha,j})|
<\frac{2|\varepsilon|E_\alpha(r;s,X_0)}{L_\alpha}. \tag{C8}
\]
For \(r=|\varepsilon|\) and fixed \(s,X_0,L_\alpha\), this is an
explicit \(O(|\varepsilon|^2)\) remainder.

\subsection{C 3 Residual roots and the exact Q08 interface}

Let \(a\) run over the distinct roots of \(f_0=F^AG\), with
multiplicity \(m_a\). A residual center \(\beta\) is a simple root of
\(G\) in the separated case, or of \(U\) in the overlap case. Choose
\(0<\rho_\beta<\min_{a\ne\beta}|\beta-a|\) and define
\[
\nu_\beta=\rho_\beta\prod_{a\ne\beta}
 (|\beta-a|-\rho_\beta)^{m_a}.
\]
The static polynomial is monic, so \(|f_0|\geq\nu_\beta\) on the
circle about \(\beta\). At fixed \(z\), the moving polynomial is even
in \(\varepsilon\). The even Cauchy tail therefore gives
\[
|H_n(z;q(\varepsilon))-f_0(z)|
\leq\mathcal B_n(Q_s,|\beta|+\rho_\beta)\frac{r^2}{s^2-r^2}
<\nu_\beta. \tag{C9}
\]
Whenever the final strict inequality holds, Rouch\'e certifies one
simple root in \(|z-\beta|<\rho_\beta\). Its analytic branch is even
and equals \(\beta+O(\varepsilon^2)\), since \(f'_0(\beta)\ne0\).

Use one common \(r,s\) in all the local tests. Give each moving cluster
the static-center radius \(R_a=rX_a\), and each residual singleton
the radius \(R_\beta=\rho_\beta\). Require also
\[
\Delta=\min_{a\ne c}(|a-c|-R_a-R_c)>0,
\qquad R=\max_a(|a|+R_a). \tag{C10}
\]
The certified roots then exhaust the degree: their counts are
\(dA+b\) in the separated case and
\((d-1)A+(A+1)+(b-1)=n\) in the overlap case. Distinct-cluster
distances are at least \(\Delta\), and all nodes have modulus at most
\(R\). For moving clusters with \(m_a>1\), put
\(\gamma_a=\delta_a-2\eta_a>0\).

For these \(n\) nodes define the square moment Vandermonde
\(V_{k+1,i}=z_i^k\), \(0\leq k<n\), \(1\leq i\leq n\).
The already accepted Q08 finite-separation envelope gives
\[
\|V^{-1}\|_2\leq\sqrt n(1+R)^{n-1}
\max_a\left\{(|\varepsilon|\gamma_a)^{-(m_a-1)}
                    \Delta^{-(n-m_a)}\right\},
\qquad
\|V\|_2\leq\sqrt{n\sum_{k=0}^{n-1}R^{2k}}.
\tag{C11}
\]
Here \(\|\cdot\|_2\) is the Euclidean operator norm. For every
singleton the first factor in braces is defined to be one; no
\(\gamma_a\) is needed. The bound follows row by row from Lagrange
interpolation: numerator coefficient \(\ell^1\) norm is at most
\((1+R)^{n-1}\), and its denominator has \(m_a-1\) within-cluster
factors and \(n-m_a\) cross-cluster factors. The inverse Frobenius
norm supplies \(\sqrt n\). The second inequality is the Frobenius
bound for \(V\).

Q08 separately already proves the fixed-degree singular-value exponents
\(\bigsqcup_a\{0,\ldots,m_a-1\}\), sorted in increasing order for
singular values in decreasing order. All comparison constants may
depend on the degree and geometry. Equations (C3)--(C10) supply verified
inputs to (C11); they are not a new singular-value-profile theorem.
Root preparation and rigorous enclosures still have a cost. Resolvent
legality, including \(1-\lambda\zeta_N^k\ne0\) for
\(\lambda=-z\), \(\zeta_N=e^{2\pi i/N}\), positive integer \(N\),
and \(0\leq k<N\), is a separate guard. A varying \(q\) here is not
a continuous family of primitive roots of one fixed order.

\subsection{C 4 A parameter only separated subfamily}

For \(b=0\) and \(A\geq2\), the Hermite directions obey
\[
\sqrt{\frac{A-1}{d}}\leq\max_j|\xi_{\alpha,j}|
\leq2\sqrt{\frac{A-1}{d}},\qquad
\frac{\pi}{\sqrt{d(A+1/2)}}\leq\delta_\alpha
\leq\frac4{\sqrt{d(A-1)}}. \tag{C12}
\]
For completeness, the Hermite recurrence realizes its zeros as the
simple real eigenvalues of the symmetric tridiagonal matrix with
off-diagonals \(\sqrt1,\ldots,\sqrt{A-1}\). Its trace square is
\(A(A-1)\), while its row-sum norm is at most \(2\sqrt{A-1}\).
The smallest gap is at most the total span divided by \(A-1\).
For the lower gap, \(u=e^{-y^2/4}\operatorname{He}_A(y)\) satisfies
\(u''+(A+1/2-y^2/4)u=0\). On any interval between consecutive zeros,
integration by parts and the Dirichlet Poincar\'e inequality give an
interval length at least \(\pi/\sqrt{A+1/2}\). Scaling by
\(1/\sqrt d\) proves (C12). The recurrence and differential equation
are classical; see NIST DLMF Tables
\href{https://dlmf.nist.gov/18.9.T1}{18.9.1} and
\href{https://dlmf.nist.gov/18.8.T1}{18.8.1}.

Put
\[
\delta_0=\frac\pi{\sqrt{d(A+1/2)}},\quad
\eta=\delta_0/4,\quad X=2\sqrt{(A-1)/d}+\eta,\quad
\mu=d^A\eta(\delta_0-\eta)^{A-1}.
\]
For fixed \(0<s<1\), let \(B=\mathcal B_{Ad}(Q_s,1+sX)\).
All \(Ad\) roots lie in the disjoint tubes if
\[
0<|\varepsilon|<
\min\left\{\frac{\mu s^{A+1}}{B+\mu s^A},
\frac{\sin(\pi/d)}X\right\}. \tag{C13}
\]
The within-cluster lower bound is
\(|\varepsilon|\delta_0/2\), and the cross-cluster lower bound is
\(2\sin(\pi/d)-2|\varepsilon|X>0\). This is (C5) with
\(|c_\alpha|=d^A\) and the exact static-center gap
\(2\sin(\pi/d)\). It provides a nonempty explicit domain for each
finite pair, not an efficient degree-uniform regime. The geometric
condition suggests small \(|\varepsilon|\sqrt{Ad}\) as parameters
vary; it is only a sufficient disk condition, not a necessary collision
threshold. The Cauchy condition may be much more restrictive.

\subsection{C 5 Exact finite examples}

Both examples use \(s=1/2\), \(r=1/10000\), \(\eta=1/8\), and
the convenient disk bound \(X=1\), which is at least each required
\(X_\alpha\). All inequalities below hold for complex
\(0<|\varepsilon|\leq r\).

\textbf{Separated example.} For \((A,d,b)=(2,3,0)\), each cluster
has \(|c|=9\), directions \(\pm\alpha/\sqrt3\), direction gap
\(>1\), and maximum direction modulus \(<3/5\). Thus
\(\mu\geq63/64\), while (C2)--(C3) give
\[
B\leq(5/4)^3(5/2)^6=\frac{1953125}{4096},\qquad
E\leq\frac{1953125}{5118976}<\frac{63}{64}.
\tag{C14}
\]
The three static centers have gap \(\sqrt3>1\); their radius-\(r\)
envelopes are disjoint. All six roots are certified, and within-cluster
distances exceed \(3|\varepsilon|/4\).

\textbf{Simple-overlap example.} For \((A,d,b)=(1,3,1)\),
\(U=g=1\) and \(\tau=3/(1-\omega^2)=1-\omega\), so
\(|\tau|=\sqrt3\). At \(-1\), \(P(x)=3x^2-\tau\) has direction
modulus \(3^{-1/4}<7/8\), gap \(>1\), and
\(\mu\geq21/64\). The other two clusters are linear, with
\(|c|=3\sqrt3\) and \(\mu>3/8\). A common deliberately coarse
majorant is \(B\leq(5/4)^4(5/2)^4=390625/4096\). Consequently
\[
E_{-1}\leq\frac{390625}{5118976}<\frac{21}{64},\qquad
E_{\mathrm{ordinary}}\leq\frac{390625}{10237952}<\frac38.
\tag{C15}
\]
All four roots are certified in disjoint static-center envelopes.

The independent finite-parameter review reconstructs these local jets
exactly in \(\mathbb Q[w]/(w^2+w+1)\), using weighted-subset Gaussian
coefficients. Its 91 Gaussian-coefficient controls through \(n=12\),
tail algebra, Hermite identities through \(A=24\), Lagrange-envelope
control with heterogeneous clusters and a singleton, and discriminant
threshold controls all pass. They support the proof, rather than prove
the general statement by extrapolation. At \(r=1/1000\) in the
separated example, the correct sufficient test is inconclusive, while
omitting \(s^{-m_\alpha}\) makes an incorrect test pass. This detects a
normalization error; it is not evidence of an actual root collision.

\subsection{C 6 Discriminant truncation error}

Set \(t=\varepsilon^2\) and
\(q(t)=\omega\exp(\operatorname{Log}(1+t)/d)\) for \(|t|<1\).
Use \((M,C)=(M_{\mathrm{sep}},C_{\mathrm{sep}})\) or
\((M_{\mathrm{ov}},C_{\mathrm{ov}})\) from Section 3, so
\(D_n(q(t))=t^M(C+O(t))\) and \(C\ne0\). For \(0<v<1\), put
\[
S=\mathcal B_n((1+v)^{1/d},1),\qquad T=n^nS^{2n-1}.
\]
On \(|t|=v\), the coefficient \(\ell^1\) norms of \(H_n\) and
its ordinary \(z\)-derivative are at most \(S\) and \(nS\).
The Sylvester matrix has \(n-1\) shifted polynomial rows and \(n\)
shifted derivative rows. Hadamard's inequality bounds its determinant
in modulus by \(T\), and hence bounds the monic discriminant by
\(T\). The normalized Cauchy tail yields, for \(0<|t|<v\),
\[
|t^{-M}D_n(q(t))-C|\leq T v^{-M}\frac{|t|}{v-|t|}. \tag{C16}
\]
For any specified relative tolerance \(\theta>0\), the strict condition
\[
|t|<\frac{\theta|C|v^{M+1}}{T+\theta|C|v^M} \tag{C17}
\]
makes this error less than \(\theta|C|\). This includes \(M=0\).
It certifies truncation error only. It does not certify accurate
floating-point evaluation of a large discriminant or infer the
individual singular spectrum from a discriminant alone.

\clearpage
\section{Appendix D Revision and source provenance}

This is a separate Version 2 working note. The original PDF and LaTeX,
the accepted all-order proof, and both finite-parameter source archives
are preserved unchanged in the companion sealed package. Sections 2--8
and Appendices A--B are unchanged. The only earlier-body edits are the
version label, the introductory description of this addition, and the
finite-parameter clarification in Section 9. Paragraph spacing is adjusted
for clean pagination. The package manifest binds
all retained sources, this revised PDF and LaTeX, and verification files.

The independent source verdict is \emph{ACCEPT} for the finite-parameter
analytic remainder, root tubes, complete separation, and discriminant
truncation bounds. It requires no mathematical correction. This
acceptance concerns its exact frozen source bytes. The present
integration is an editorial revision; independent document-fidelity
review of it is a separate step. No historical-priority, implementation,
or research-ranking upgrade follows from this addition.

\textbf{Original working note.} PDF Library identity
\nolinkurl{libfile_72ddfda84adc8191b5ed994749139490}; LaTeX identity
\nolinkurl{libfile_46168cfa2b3481918766ba14fbc98d66}, both version 0.
Their SHA256 values, respectively, are

\nolinkurl{c65f50a5dc3acc0580f35e17ead1adbd9b5dc6a7d7f42d4970e8e1c3809e9372}

\nolinkurl{23462defbe98daf588869a27e46f4ad94f67fd02a09e0196d31bf231aa15c1f0}

\textbf{Accepted finite-parameter source.} Archive Library identity
\nolinkurl{libfile_730b9b6c2d708191bc51cf5419553afe}, version 0.
Archive SHA256:

\nolinkurl{2f2180470328e69db523339981b83918698b62f2589ad750010806b4fc8923ef}

The accepted LaTeX member SHA256 is

\nolinkurl{3b472a7adb4a157c52c6ad4f8da5e1d94717e360bce8326e8c4d8732cfbcc100}

\textbf{Independent finite-parameter review.} Archive Library identity
\nolinkurl{libfile_30c794c5011c819196acc7d3104d8aad}, version 0.
Archive SHA256:

\nolinkurl{b5a00e74f00f1f9d2b5ebfc1afb9479bb950efca792504d63a1c79ce6ffc3e72}

The accepted review member SHA256 is

\nolinkurl{a412af68fbeffe1d6f332280a5b59e3ac74d91481c492232981319c51e0b6033}

\textbf{Q08 interface dependencies.} The source \emph{Structured fusion
second pass}, Section S3, supplies the fixed-degree singular profile
and finite inverse envelope. Its SHA256 is

\nolinkurl{02ec40df49f0585c57f98d3af7cbb9867216773f98a2998db9d51237fd69b787}

The separate resolvent and root-path warning is in \emph{Fusion proofs},
Theorem T3, SHA256

\nolinkurl{4afd27bacea1bd896d8b44e7e4c1eb51c1b24b3db997ceca9e554e69f81202bb}

The accepted all-order proof and review are pinned in Appendix B.
Unverified approximate root lists are not certificates: numerical use
of Appendix C requires exact algebraic data or rigorous outward-rounded
enclosures. The historical-source limitations in Section 9 remain in
force.

\end{document}
